% Generated draft: requires mathematical and source review.
\documentclass[11pt]{article}
\usepackage[margin=25mm]{geometry}
\usepackage{fontspec}
\setmainfont{Latin Modern Roman}
\usepackage{amsmath,amssymb,mathtools,unicode-math}
\setmathfont{Latin Modern Math}
\usepackage{enumitem,needspace,xurl,hyperref}
\setlength{\parindent}{0pt}
\setlength{\parskip}{5pt}
\setlength{\emergencystretch}{3em}
\newcommand{\N}{\mathbb{N}}
\newcommand{\Z}{\mathbb{Z}}
\newcommand{\Q}{\mathbb{Q}}
\newcommand{\R}{\mathbb{R}}
\newcommand{\C}{\mathbb{C}}
\newcommand{\F}{\mathbb{F}}
\newcommand{\A}{\mathbb{A}}
\newcommand{\PP}{\mathbb{P}}
\newcommand{\E}{\mathbb{E}}
\newcommand{\eps}{\varepsilon}
\newcommand{\dd}{\,\mathrm{d}}
\begin{document}
\section*{UnsolvedMath ID: 30002444}
\textbf{Source problem code:} catalogue entry
\subsection{Definitions and mathematical statement}
The supplied record asks two homogeneity questions. The formulation below uses the customary complex-geometric interpretation of a Fano manifold; the record does not explicitly specify the base field. Under this interpretation, let \(X\) be a smooth connected complex projective variety with ample anticanonical line bundle \(-K_X\), and assume
\[
 b_2(X):=\dim_{\Q}H^2(X,\Q)=1.
\]
Choose a covering family \(\mathcal K\) of minimal rational curves on \(X\). For a general point \(x\in X\), its variety of minimal rational tangents is
\[
 \mathcal C_x:=\overline{\{[T_xC]\in\PP(T_xX): C\in\mathcal K,\ x\in C,\ C\text{ is smooth at }x\}},
\]
where the bar denotes Zariski closure and \(\PP(T_xX)\) parametrizes one-dimensional tangent subspaces. Here ``general'' means that the condition holds on a nonempty Zariski open subset of \(X\). The record does not identify \(\mathcal K\), specify the precise minimality convention, or say how the assertion is quantified over possible choices of this family.

A projective subvariety is called nonlinear here if it is not a projective linear subspace. Call \(X\) homogeneous if its algebraic automorphism group \(\operatorname{Aut}(X)\) acts transitively on \(X\), and quasi-homogeneous if this group has a Zariski open dense orbit.

For clarity about the Legendrian hypothesis, a contact projective space can be described by a complex vector space \(V\) of dimension \(2m+2\) with a nondegenerate alternating bilinear form \(\omega\). Its contact distribution at a line \(\ell\subset V\) is
\[
 D_{\ell}=\operatorname{Hom}(\ell,\ell^{\perp_\omega}/\ell)
 \subset T_{\ell}\PP(V).
\]
A smooth subvariety \(L\subset\PP(V)\) is Legendrian if \(\dim L=m\) and \(T_{\ell}L\subset D_{\ell}\) for every \(\ell\in L\).

\textbf{Legendrian question.} For every \(X\) as above, suppose that its VMRT at a general point is a connected nonlinear Legendrian submanifold of the contact projectivization intended in the source. Must \(X\) be homogeneous? The separately stated weaker question asks whether \(X\) must at least be quasi-homogeneous. The record does not identify the vector space underlying that contact projectivization, its contact structure, or the embedding of \(\mathcal C_x\) into it. In particular, it does not specify whether the ambient space is \(\PP(T_xX)\) or the projectivization of a proper tangent subspace. Thus this ambient hypothesis remains unresolved; no contact structure on \(X\) is added here.

\textbf{Hypersurface question.} For every \(X\) as above, suppose that \(\mathcal C_x\subset\PP(T_xX)\) is a smooth nonlinear hypersurface for general \(x\). Must \(X\) be homogeneous? The weaker question asks whether \(X\) must be quasi-homogeneous. In this usual ambient interpretation, the hypersurface has codimension one and degree at least two. Unlike the Legendrian hypothesis, this hypothesis does not explicitly require connectedness. The supplied wording does not explicitly identify the hypersurface ambient space, so the use of the natural VMRT ambient \(\PP(T_xX)\) requires confirmation.

The supplied dated assessment labels the overall record partially solved, reports that the weaker hypersurface version was settled, and describes the connected nonlinear Legendrian version as remaining open. It gives no precise settling theorem or separate citation. These are supplied status claims, not independently verified mathematical conclusions.
\subsection{Short English statement}
The questions ask whether a Fano manifold with \(b_2=1\) must be homogeneous, or at least have an open automorphism orbit, when its general variety of minimal rational tangents is either connected, nonlinear and Legendrian, or a smooth nonlinear hypersurface. The minimal-curve family and some ambient contact data are unspecified. The supplied dated curation reports partial resolution; this draft does not independently verify that assessment.
\subsection{Sources}
\begin{itemize}
\item Supplied local UnsolvedMath registry record, permanent ID \(30002444\), title: \emph{Homogeneity from Legendrian Varieties of Minimal Rational Tangents}; exact source problem code: \texttt{catalogue entry}. Collection: Oberwolfach Reports Open Problems. No individual proposer is identified; the supplied proposed year is 2013.
\item Supplied source citation: \emph{Mini-Workshop: Quaternion K\"ahler Structures in Riemannian and Algebraic Geometry} (2014), Oberwolfach Reports, pp.~3115--3145, DOI \texttt{10.4171/owr/2013/53}. Supplied source URL and literature-source link, labelled ``original OWR report'': \url{https://doi.org/10.4171/owr/2013/53}. The record describes these reports as documenting mathematical workshops and their open-problem discussions; no problem-specific page or named contributor is supplied.
\item Literature checked at: \texttt{2026-08-21}. This date and the assessment are supplied dated curation, not independently verified. The supplied assessment states that the weaker hypersurface version was settled while the connected nonlinear Legendrian version remains open. The record itself requests independent verification of time-sensitive status and cited claims. Its sole supplied literature reference is the original report above; no separate citation for the reported resolution is provided.
\end{itemize}
\end{document}
